% [原文 p.57]
\begin{enumerate}[label=（\arabic*）,leftmargin=2.6em,itemsep=1pt,start=2]
\item （极小模型）若 $K_{X_i}$ 为 nef，则同样定义 $X^*:=X_i$。在试图更好地理解 $X^*$ 的过程中，可以期待更有效地利用标准类为 nef 这一事实。
\end{enumerate}

上述纲领看起来更像是并不完整、只是罗列希望的东西。第一，它只在二维与三维完成。更高维时，人们尚不知道步骤 2 (2) 的翻转是否总存在。纲领最终是否必定停止也不清楚，即翻转有可能无限次重复。最后，所得结果（步骤 3）究竟是否有用也不清楚。诚然，簇的某种全局性质改善了，但代价是产生了奇点。

开发必要的技术花了约 10 年，但现在已查明：高维代数几何中许多结果的证明，其实质步骤正是由上述纲领给出的。\hfill$\square$

% [Tutor Note: 上述四点保留是纲领的「信任状」：它在二维、三维成立，且大量高维定理的证明实质上都依赖它。因此后续各章的重点不是怀疑纲领，而是补上缺失的技术环节。]

\begin{example}{2.15}
这里给出若干高维极射线收缩的例子。
\begin{enumerate}[label=（\arabic*）,leftmargin=2.6em,itemsep=1pt]
\item 设 $Y$ 为非奇异射影簇，$Z\subset Y$ 为非奇异闭子簇。此时爆破的逆 $B_ZY\to Y$ 是极射线收缩态射。
\item 令
\[
X=\PP_{\PP^n}\bigl(\cO_{\PP^n}(1)\oplus\cdots\oplus\cO_{\PP^n}(1)\oplus\cO_{\PP^n}\bigr),
\]
其中共有 $k$ 个直和项 $\cO_{\PP^n}(1)$。此时直和成分 $\cO_{\PP^n}$ 给出自然嵌入 $\PP^n\cong E\subset X$。
若 $k\le n$，则直线 $L\subset E\subset X$ 生成 $\NE(X)$ 的 $K_X$-负极射线。相应的收缩把 $E$ 压成一点，在 $E$ 之外为同构。因此当 $k\ge2$ 时例外集不是除子；当 $\dim X=n+k\ge4$ 时，用此法可以造出这样的例子。
\item 设 $Y$ 为非零线性映射 $\C^{n+1}\to\C^n$ 的全体按标量倍数取商所得的空间。即对 $f,g\colon\C^{n+1}\to\C^n$，$f\sim g$ 意为：存在标量 $\lambda\ne0$，使 $f(x)=g(\lambda x)$ 对一切 $x\in\C^{n+1}$ 成立。因 $Y\cong\PP^{n(n+1)-1}$，故 $Y$ 非奇异。$X$ 是 $(g,L)$
% [原文 p.58]
的空间，其中 $g\in Y$，$L$ 是含于 $g$ 的核空间中的 $1$ 维线性子空间。
设 $f\colon X\to Y$ 为自然态射，则可证 $f$ 是极射线收缩。由 $p(g,L)=L$ 给出 $X$ 到 $\PP^n$（即 $\C^{n+1}$ 的 $1$ 维子空间全体）的自然态射 $p$。$p$ 为满射，其纤维都是 $(n^2-1)$ 维射影空间，$X$ 也仍非奇异。
再定义
\[
F=\{g:\operatorname{rank}g\le n-1\},\qquad E=\{(g,L):\operatorname{rank}g\le n-1\}.
\]
考虑 $p$ 在 $E$ 上的限制，则 $E$ 成为 $\PP^n$ 上的纤维丛。不过 $L$ 上的纤维，是退化线性映射 $\C^{n+1}/L\to\C^n$ 全体的射影化。因此 $E$ 不可约。若 $g\in F$，则 $f^{-1}(g)$ 是 $(n-\operatorname{rank}g)$ 维射影空间。故对一般的 $g\in F$，它是 $\PP^1$。若 $n>2$，则存在 $g\in F$ 使 $\operatorname{rank}g=n-2$，于是 $f^{-1}(g)\cong\PP^2$。
这意味着 $f$ 不可能是非奇异爆破。事实上可以证明：$F$ 在 $g$ 处有奇点，当且仅当 $\operatorname{rank}g\le n-2$。
\item 设 $C$ 为非奇异射影曲线，$p\in C$。记 $S^mC$ 为 $C$ 的 $m$ 次对称积，$J(C)$ 为 $C$ 的雅可比（Jacobi）簇。注意 $S^mC$ 是非奇异射影簇。对每个 $m$，点 $p$ 给出态射 $u_m\colon S^mC\to J(C)$，它在 $m\ge g(C)$ 时为满射，在 $m=g(C)$ 时为双有理的。当 $m\ge g(C)$ 时，$u_m$ 是 $S^mC$ 的 $K$-负极射线收缩（参见 [Kol96, III.1.6.6]）。
这给出非奇异簇之间极射线收缩且例外集复杂的大量例子。
\end{enumerate}
\end{example}

% [Tutor Note: 例 2.15(3) 的纤维从 $\PP^1$ 跳到 $\PP^2$，说明极射线收缩不必来自爆破。这种「小收缩」正是必须在容许奇点的簇类中工作的原因，也是翻转出现的场合。]

\section{极小模型纲领的推广}
\takeaway{本节介绍极小模型纲领的 4 个有益推广：相对化、解析的情形、带群作用的簇与对数簇。它们共用同一套锥定理与收缩定理，因此纲领的框架可以整体照搬。}

极小模型纲领的 4 个有益推广如下：
\begin{enumerate}[label=（\arabic*）,leftmargin=2.6em,itemsep=1pt]
\item 相对化
\item 解析的情形
% [原文 p.59]
\item 带群作用的簇
\item 对数（对数的，log）簇
\end{enumerate}

\begin{example}[相对化]{2.16}
设 $f\colon X\to Y$ 为射影态射。为了引入 $\R$-向量空间 $N(X/Y)$，取一切满足 $f(C)$ 为一点的不可约曲线 $C\subset X$ 为生成元，并按
\[
\sum a_iC_i\sim0\iff\sum a_i(C_i\cdot D)=0
\]
定义等价关系，其中 $D$ 取遍 $X$ 上的一切 Cartier 除子。$N_1(X/Y)$ 的元素称为相对 $1$-循环（relative $1$-cycle）。$N_1(X/Y)$ 的维数称为 $X/Y$ 的相对 Picard 数（relative Picard number），记作 $\rho(X/Y)$。$\NE(X/Y)\subset N(X/Y)$ 定义为由被 $f$ 压成一点的有效曲线类生成的凸锥。$\overline{NE}(X/Y)$、$\overline{NE}(X/Y)_{D\ge0}$ 等相对概念也与绝对情形同样定义。
若 $X$ 与 $Y$ 固有，则交点数定义双线性形式
\[
N_1(X/Y)\times\bigl(NS(X)/f^*NS(Y)\bigr)\otimes_{\Z}\R\to\R.
\]
一般地，$N_1(X/Y)$ 与 $(NS(X)/f^*NS(Y))\otimes_{\Z}\R$ 连维数\footnote{两者都是 $\rho(X/Y)$ 的自然候选，但本书选择了着眼于曲线的定义。不过如下文所述，在本书所处理的状况下，多数场合它们一致。}都不同，故不能说该双线性形式非退化（参见例 1.46）。不过，$f$ 为双有理的且 $Y$ 为 $\Q$-分解的时，容易看出它非退化。
设 $f\colon X\to Y$ 与 $g\colon Y\to Z$ 为固有态射，$f$ 为满射且纤维连通。此时得到复形（complex）
\[
0\to N_1(X/Y)\to N_1(X/Z)\to N_1(Y/Z)\to 0.
\]
它一般不是正合列，但若下列之一成立，则为正合列。
\begin{enumerate}[label=（\arabic*）,leftmargin=2.6em,itemsep=1pt]
\item $(X,\Delta)$ 为 klt（参见后文定义 2.34），且对适当的有效 $\Q$-除子 $\Delta$，$-K_X-\Delta$ 为 $f$-丰富 [KMM87, 3-2-5]；
\item $X$ 与 $Y$ 只有有理奇点，且 $R^1f_*\cO_X=0$ [KM92, 12.1.5]；
\item $f$ 为双有理的，且 $Y$ 为 $\Q$-分解的（利用 [KM92, 12.1.5] 的论证）。
\end{enumerate}
锥定理与收缩定理（参见 §3.1）与绝对情形相同（证明也相同）。
% [原文 p.60]
在证明非奇异情形的锥定理（定理 1.24）所用的论证中，应注意：若出发的曲线 $C$ 满足 $f(C)$ 为一点，则由它得到的一切曲线都被 $f$ 送到同一点。
设上述 $X$ 为三维簇且非奇异（或为 $\Q$-分解的且只有终端奇点）。若利用 [Mor88] 所建立的三维翻转的存在性，则不断进行 $Y$ 上的收缩或翻转，可到达满足下列之一的、$\Q$-分解的且只有终端奇点的态射 $f'\colon X'\to Y$。
\begin{enumerate}[label=（\arabic*）,leftmargin=2.6em,itemsep=1pt]
\item $X'$ 为 $Y$ 上的极小模型，即 $K_{X'}$ 为 $f'$-nef；或
\item $X'$ 具有 $Y$ 上的 $\Q$-Fano 纤维结构 $g'$，即有如下交换图式。
\end{enumerate}
\begin{center}
\begin{tikzcd}[column sep=5em, row sep=3em]
X' \arrow[rr, "g'"] \arrow[dr, "f'"'] & & Z' \arrow[dl, "h'"] \\
{} & Y & {}
\end{tikzcd}
\end{center}
其中 $g'$ 具有连通纤维，$\dim Z'<\dim X'$，且 $-K_{X'}$ 为 $g'$-丰富。
证明中与绝对情形不同的只是技术性细节。详见 [KMM87]。
\end{example}

\begin{example}[解析的情形]{2.17}
首先，有 3 个必须在用语上特别注意之处。
对解析流形而言，即使紧的情形，按记号 0.4 的意义也可能不存在「标准除子」。不过本书只用到它的线性等价类，即对偶化层（dualizing sheaf）与标准丛。因此这纯粹是记号滥用，不会引起任何问题。
在相对的情形，有时要用生成纤维的概念。这在解析情形也没有定义，但若改用充分一般的纤维，就不会有任何问题。
最后是 $\Q$-分解性的概念。如前所述，解析空间 $X$
% [原文 p.61]
为 $\Q$-分解的（或与局部解析的 $\Q$-分解性相对，称全局解析的 $\Q$-分解的）是指：$X$ 上的一切 Weil 除子都是 $\Q$-Cartier。$X$ 为局部解析 $\Q$-分解的是指：$X$ 的一切 Euclid 开集都是全局解析 $\Q$-分解的。
对这些概念必须十分小心。原因在于：代数（或全局解析）的 $\Q$-分解性并不蕴含全局解析（或局部解析）的 $\Q$-分解性。（$X^{an}$ 上的 Cartier 除子未必是代数的；而在解析情形，开集上的 Weil 除子也未必能延拓为整个空间上的 Weil 除子。）若射影簇 $X$ 为 $\Q$-分解的，则伴随的解析空间 $X^{an}$ 是全局解析 $\Q$-分解的。这由 Chow 定理保证：$X^{an}$ 的一切全局除子都是代数的（参见 [Har77, App.B]）。
考虑到这些，已知解析空间版 MMP 的情形，是对某个满足有限性条件的解析空间 $Y$ 的射影态射 $f\colon X\to Y$ 而言的。（关于代数簇的）相对代数版也同样成立，因为所需的相对消没定理在这种情形也成立。详见 [Nak87]。
$f$ 不为射影时，所知甚少。[Har77, p.443] 所介绍的广中的例子表明：定理 1.32 对非奇异紧三维复流形不成立。详见 7.80。
\end{example}

% [Tutor Note: 解析情形没有「印在纸上」的标准除子，只能用标准丛。这提醒我们：定义 2.22 中的除子本质上是线丛的数据，而不是空间中的真实除子。]

\begin{examplep}[带群作用的簇]{2.18}
设射影簇 $X$ 非奇异，或为 $\Q$-分解的且只有终端奇点，并设有限群 $G$ 作用于其上。此时对 $\NE(X)^G\subset N(X)^G$ 有锥定理与收缩定理。这里 $M^G$ 表示 $M$ 的 $G$-不变元全体。不同之处在于：$K_X$ 是 $G$-不变的，故 $G$-极射线对应于通常意义下极射线的 $G$-轨道，而它对应于极锥面。因此收缩定理也意味着 $G$-不变极锥面的收缩。
在其它设定中也有应用。例如，设 $X$ 是定义在域 $k$ 上的曲面。把 $k$ 上的极小模型按 $G=\operatorname{Gal}(\bar k/k)$ 取，即可得到 $G$-极小模型。其中 $\bar k$ 是 $k$ 的代数闭包。$G$ 不是有限群，但它对 $X_{\bar k}$ 的 Néron--Severi 群的作用通过有限商进行，故 $G$-极小模型的构
% [原文 p.62]
成与代数闭域上的情形同样进行。
若 $X$ 是带有限群 $G$ 作用的非奇异复射影曲面，则除细微差别外，分类与无作用的情形相同。$G$-极射线由形如 $C=\sum C_i$ 的 $1$-循环生成，其中 $C_i$ 是被 $G$ 互相移动的不可约有理曲线。
(1) 若 $(C^2)<0$，则容易看出 $C_i$ 是互不相交的 $(-1)$-曲线。因此可把一切 $C_i$ 压成非奇异点。
(2) 若 $(C^2)=0$，则 $C$ 的任意连通分支或为不可约曲线，或为只在一点横交的两条 $(-1)$-曲线。此时收缩把 $X$ 变成非奇异曲线上的二次曲线丛。
(3) 若 $(C^2)>0$，则 $N(X)^G=\R$ 且 $-K_X$ 为丰富。特别地，$X$ 成为 Del Pezzo 曲面。
关于三维簇，在带群作用的奇点消解之后，由带群作用的极小模型纲领可得如下结果。（$GQ$-分解的是指任意 $G$-不变 Weil 除子都是 $\Q$-Cartier。）
\end{examplep}

\begin{theorem*}{2.19}
任意固有的三维 $G$-簇 $X$，都与某个 $GQ$-分解的、且只有终端奇点的 $G$-簇 $Y$ 连同 $G$-作用双有理等价，且 $Y$ 满足下列之一。
\begin{enumerate}[label=（\arabic*）,leftmargin=2.6em,itemsep=1pt]
\item $K_Y$ 为 nef；或
\item 存在 $G$-态射 $f\colon Y\to Z$，使 $-K_Y$ 为 $f$-丰富，且 $\dim Z<\dim X$。
\end{enumerate}
\end{theorem*}

\begin{example}[对数簇]{2.20}
在极小模型纲领发展的过程中，人们逐渐明白：不仅标准类，其变种也值得作为研究对象。这一方向的研究最初的契机，是饭高关于开簇的研究 [Iit77]。
代替 $K_X$，考虑正规簇 $X$ 上形如 $K_X+D$ 的除子。其中 $D=\sum d_iD_i$ 是 $D_i$ 互不相同、系数 $d_i$ 满足条件 $0\le d_i\le1$ 的形式线性组合。理由至少有 4 个。
\begin{enumerate}[label=（\arabic*）,leftmargin=2.6em,itemsep=1pt]
\item 柔性：通过适当选取 $D$，在 $K_X$ 很小（例如 $K_X\equiv0$）
% [原文 p.63]
时，或在 $K_X$ 不是 $\Q$-Cartier 时，也能进行研究。
\item 归纳证明。近年来，人们考察了几种把极小模型的某些问题归约到更低维问题来处理的方法。多数情况下，只有当处理的是对 $(X,D)$、即 $K_X+D$ 时才可能归约。这类情形中的若干已在 [K+92] 中处理。
\item 开簇。这是 [Iit77] 的想法。设 $X$ 为非奇异簇，$X\subset\bar X$ 为非奇异紧化，且 $D=\bar X-X$ 为正规交叉除子。
不可思议的是，用沿 $D$ 具有对数型极点的正则微分形式所构成的许多向量丛，其上同调群不由紧化 $\bar X$ 决定，而只由 $X$ 决定。这样的例子有
\[
H^j\bigl(\bar X,\Omega^j_{\bar X}(\log D)\bigr),\qquad H^0\bigl(\bar X,\cO(m(K_{\bar X}+D))\bigr)\quad(m\ge0)
\]
等。因此，若要研究 $X$ 的性质，考察除子 $K_{\bar X}+D$ 就是自然的。
\item 纤维空间。以椭圆曲面上的小平标准除子公式为例（参见例如 [BPdV84, V.12.1]）。设 $f\colon S\to C$ 为极小椭圆曲面，$m_iF_i=f^*(c_i)$ 为多重纤维。此时下式即标准除子公式。
\[
K_S=f^*K_C+f^*(f_*K_{S/C})+\sum(m_i-1)F_i
\equiv f^*\Bigl[K_C+(f_*K_{S/C})+\sum\Bigl(1-\frac{1}{m_i}\Bigr)[c_i]\Bigr].
\]
这里 $f_*K_{S/C}$ 表示可逆层 $f_*\cO_S(K_S-f^*K_C)$ 所定的除子类。因此研究 $K_S$，就归结为对具有有理系数的 $D$ 研究形如 $K_C+D$ 的除子。同样的事情对一般纤维具有平凡标准除子的纤维空间也成立。
\end{enumerate}
不过，要使证明行得通，对 $(X,D)$ 必须加上某种技术性条件。本质在于：不允许除子 $K+D$ 过于奇异。这一点的确切含义将在下节阐明。
\end{example}

% [Tutor Note: 对数簇的 $D$ 可看作「允许的边界」：系数越小，允许的奇点越坏。把 $D$ 加到饱和后，$K_X+D$ 在边界上的行为就决定了容许的奇点类。]

% [原文 p.64]
\begin{example}[相对化、解析的情形、带群作用的簇、对数簇]{2.21}
研究上述推广中若干种同时发生的情形也是可能的。例如，可以研究相对的、带 $G$-作用的对数（对数的）极小模型纲领。一切结果都能推广到这一设定。
\end{example}

\section{极小模型纲领中的奇点}
\takeaway{要继续纲领，必须理解收缩极射线时产生的奇点。差异系数 $a(E,X)$ 是衡量簇奇异程度的尺度，它只由除子 $E$ 对应的赋值决定，与所选模型无关。}

如定理 1.32 所示，收缩极射线会产生奇异的簇。要继续纲领，就必须理解过程中产生的奇点。为此需要衡量簇「奇异程度」的「尺度」。这一新尺度称为差异系数。

\begin{definition}{2.22}
设 $X$ 为正规簇，且存在整数 $m>0$ 使 $mK_X$ 为 Cartier。用 $f\colon Y\to X$ 表示来自正规簇 $Y$ 的双有理态射（不必固有）。设 $E\subset Y$ 为不可约例外除子，$e\in E$ 为 $E$ 的一般点，$\{y_i\}$ 为 $e\in Y$ 处的局部坐标系，且 $E=(y_1=0)$。此时在 $e$ 附近，$f^*$ 把 $\cO_X(mK_X)$ 的局部生成元映为
\[
y_1^{m\cdot a(E,X)}\,u\,(\rd y_1\wedge\cdots\wedge \rd y_n)^{\otimes m},
\]
其中 $u$ 为单位。这里 $a(E,X)$ 是使 $m\cdot a(E,X)$ 为整数的有理数。$a(E,X)$ 称为 $E$ 关于 $X$ 的差异系数（discrepancy）。当然，$a(E,X)$ 与 $m$ 的取法无关。
若 $f\colon Y\to X$ 为固有双有理态射，且 $K_Y$ 为 Cartier（例如 $Y$ 非奇异），则 $mK_Y$ 线性等价于
\[
f^*(mK_X)+\sum_i\bigl(m\cdot a(E_i,X)\bigr)E_i,
\]
其中 $E_i$ 取遍一切例外素除子。若改用数值等价，除以 $m$ 后即可写成
\[
K_Y\equiv f^*K_X+\sum_i a(E_i,X)E_i.
\]
\noindent
\end{definition}

% [Tutor Note: 差异系数的符号是奇点分类的标尺：对一切 $E$ 都有 $a(E,X)>0$ 即终端奇点，$a(E,X)\ge0$ 即标准奇点。定义 2.28 的 $\operatorname{discrep}$ 正是把这种「最小符号」量化。]

\begin{remark}{2.23}
设 $k(X)$ 为 $X$ 的有理函数体。对素除子 $E$，$\cO_{E,X}\subset k(X)$
% [原文 p.65]
表示 $E$ 的泛点 $e$ 处 $X$ 的局部环 $\cO_{e,X}$，而它是 $k(X)$ 的与赋值 $v(E,X)$（下文如此记）对应的离散赋值环。$k(X)$ 的这种赋值称为代数赋值（algebraic valuation）。代数赋值的抽象刻画由引理 2.45 给出。
设 $f\colon Y'\to X$ 为另一双有理态射，$E'\subset Y'$ 为满足 $v(E,Y)=v(E',Y')$ 的素除子。这等价于：$Y\to X\leftarrow Y'$ 在 $E$ 与 $E'$ 的泛点 $e,e'$ 处都同构。此时由定义容易看出 $a(E,X)=a(E',X)$。
这样，$a(E,X)$ 只由赋值 $v(E,X)$ 决定，与 $f$ 及 $Y$ 的取法无关。正因如此，$f$ 与 $Y$ 在记号中被省略。
\end{remark}

\begin{definition}{2.24}
设 $X$ 为簇。$f\colon Y\to X$ 为（不必固有的）双有理态射，$Y$ 正规，$E\subset Y$ 为素除子。这样的 $E$ 称为 $X$ 上空的除子（divisor over $X$）；或者把它与其赋值等同，称为 $X$ 上的代数赋值。$f(E)\subset X$ 的闭包称为 $E$ 在 $X$ 上的中心（center），记作 $\operatorname{center}_XE$。如上所见，中心只由赋值 $v(E,X)$ 决定。当 $E$ 为 $f$-例外时，称赋值 $v(E,X)$ 为 $X$ 上例外的。
\end{definition}

定义 2.22 可以推广到使 $K_X+\Delta$ 为 $\Q$-Cartier 的对 $(X,\Delta)$。

\begin{definition}{2.25}
$(X,\Delta)$ 是正规簇 $X$ 与 $\Q$-除子 $\Delta$ 的对。这里用互不相同的素除子 $D_i$ 把 $\Delta$ 表示为 $\Delta=\sum a_iD_i$（系数 $a_i$ 可为任意有理数）。另外，假设存在整数 $m>0$ 使 $m(K_X+\Delta)$ 为 Cartier 除子。设 $f\colon Y\to X$ 为来自正规簇 $Y$ 的双有理态射。用 $E\subset Y$ 表示 $f$ 的例外集，$E_i\subset E$ 取遍一切不可约例外除子。又用
\[
f_*^{-1}\Delta:=\sum a_i f_*^{-1}D_i
\]
表示 $\Delta$ 的双有理变换。线丛 $\cO_Y(m(K_Y+f_*^{-1}\Delta))|_{(Y-E)}$ 与 $f^*\cO_X(m(K_X+\Delta))|_{(Y-E)}$ 自然同构。因此存在有理数 $a(E_i,X,\Delta)$，使 $m\cdot a(E_i,X,\Delta)$ 为整数，且
\[
\cO_Y\bigl(m(K_Y+f_*^{-1}\Delta)\bigr)\cong f^*\cO_X\bigl(m(K_X+\Delta)\bigr)\Bigl(\sum_i m\cdot a(E_i,X,\Delta)E_i\Bigr),
\]
% [原文 p.66]
其中右端的括号表示按该除子作扭转。由定义，$a(D_i,X,\Delta)=-a_i$；而对与任何 $D_i$ 都不同的任意素除子 $D\subset X$，有 $a(D,X,\Delta)=0$。这个 $a(E,X,\Delta)$ 称为 $E$ 关于 $(X,\Delta)$ 的差异系数（discrepancy）。不致产生混乱时，也简写为 $a(E)$。
与 $\Delta=0$ 的情形一样，$a(E_i,X,\Delta)$ 只由 $E_i$ 决定，与 $f$ 无关。
\end{definition}

\begin{notation}{2.26}
用数值等价，除以 $m$ 后可写成
\[
K_Y+f_*^{-1}\Delta\equiv f^*(K_X+\Delta)+\sum a(E_i,X,\Delta)E_i,
\]
其中和式取遍一切例外的 $E_i$；又可写成
\[
K_Y\equiv f^*(K_X+\Delta)+\sum a(E_i,X,\Delta)E_i,
\]
其中和式取遍一切 $E_i$。后文常说「$K_Y\equiv f^*(K_X+\Delta)+A$」并引用这些标准丛公式。遇到这种情形，$A$ 总按上述方式选取：即把非例外除子的系数取成期望的值，也就是选取 $A$ 使 $f_*A=-\Delta$。\hfill$\square$
\end{notation}

\begin{lemma}{2.27}
记号同上。设 $\Delta'$ 为有效 $\Q$-Cartier 除子。此时对 $X$ 上空的一切除子 $E$，都有 $a(E,X,\Delta)\ge a(E,X,\Delta+\Delta')$。进一步，取「$>$」当且仅当 $\operatorname{center}_XE\subset\operatorname{Supp}\Delta'$。
\end{lemma}

% [Tutor Note: 引理 2.27 的单调性说明：把边界 $\Delta$ 换成更大的 $\Delta+\Delta'$，差异系数只会变小。边界越大，标准类越「退化」，允许的奇点也越坏，这与 $K_X+D$ 的几何直觉一致。]

本书的观点中最重要的是：当 $E$ 取遍各种除子时，$a(E,X,\Delta)$ 所取的最小值。下面定义其中几个。

\par\medskip\noindent\textbf{定义 2.28}\quad $(X,\Delta)$ 的差异系数（discrepancy）由下式给出：
\[
\operatorname{discrep}(X,\Delta):=\inf\{a(E,X,\Delta)\},
\]
其中 $E$ 取遍 $X$ 上空的一切例外除子（即取遍一切双有理态射 $f\colon Y\to X$ 的一切例外除子）。等价地说，$\operatorname{discrep}(X,\Delta):=\inf\{a(v,X,\Delta)\}$，其中 $v$ 取遍 $X$ 上一切例外的代数赋值。
另外，$(X,\Delta)$ 的全差异（total discrepancy）由下式定义。
